\documentclass[a4paper, landscape, 10pt, twocolumn]{article}
\usepackage{xeCJK}
\usepackage[left=1.5cm, right=1.5cm, top=1.5cm, bottom=1.5cm]{geometry}
\usepackage{tabularx}
\usepackage{amsmath}
\usepackage{fancyhdr}
\usepackage[table]{xcolor}
\pagestyle{fancy}
\lhead{高中物理}
\rhead{南昌学而思}
\cfoot{generated by Linghui Ding}
\usepackage{multicol}
\setlength{\columnseprule}{0.4pt}
\begin{document}
\rowcolors{2}{white}{blue!10}
\renewcommand{\arraystretch}{1.5}
\noindent
\begin{tabularx}{\columnwidth}{XX}
	\rowcolor{blue!30}
	\multicolumn{2}{l}{\parbox[c][2.5\ht\strutbox][c]{0.9\columnwidth}{\Large 静电场}} \\
    场强定义&$E=\frac Fq$\\
    电势定义&$\varphi = \frac Wq$\\
    电势差定义&$U_{AB}=\varphi_A-\varphi_B$\\
    电场力做功&$W=qU_{AB}=W_A-W_B=\varSigma qE\Delta d$\\
    \multicolumn{2}{l}{\textbf{匀强电场规律}}\\
    电势差与场强关系&$U=Ed$\\
	\multicolumn{2}{l}{\textbf{点电荷}}\\
    库仑定律&$F=\frac{kQq}{r^2}$\\
    点电荷电场强度&$E=\frac{kQ}{r^2}$\\
    点电荷电势&$\varphi=\frac{kQ}{r}$\\
\end{tabularx}
\textbf{注意}\\
\begin{tabularx}{\columnwidth}{XX}
	\rowcolor{blue!30}
	\multicolumn{2}{l}{\parbox[c][2.5\ht\strutbox][c]{0.9\columnwidth}{\Large 电路基本元件}} \\
    电流定义式&$I=\frac Qt$\\
    电流决定式&$I = nqSv$\\
    电阻定义式&$R=\frac UI$\\
    电阻决定式&$R=\rho\frac lS$\\
    电容定义式&$C=\frac QU$\\
    电容决定式&$C=\frac{\epsilon S}{4\pi kd}$\\
    电源电动势&$E=\frac{W_{\text{非静电力}}}{q}$\\
\end{tabularx}
\textbf{注意}\\
\begin{tabularx}{\columnwidth}{XX}
	\rowcolor{blue!30}
	\multicolumn{2}{l}{\parbox[c][2.5\ht\strutbox][c]{0.9\columnwidth}{\Large 闭合电路欧姆定律}} \\
    不含外电路电压&$E=I(R+r)$\\
    不含内阻和电动势&$U=IR$\\
    不含外电阻&$U=E-Ir$\\
    不含电流&$U=E\frac{R}{R+r}$\\
    串联规律&$R=R_1+R_2+\cdots R_n$\\
    并联规律&$\frac1R=\frac1{R_1}+\frac1R_2+\cdots \frac1R_n$\\
    \multicolumn{2}{l}{\parbox[c][2.5\ht\strutbox][c]{0.9\columnwidth}{\Large 电表改装}} \\
    电流表简易改装&$R=\frac{I_gR_g}{I_{\text{量程}}-I_g}\approx \frac{I_gR_g}{I_{\text{量程}}}$\\
    电压表简易改装&$R=\frac{U_{\text{量程}}}{I_g}-R_g$\\
    欧姆表改装&$E=i(r+R_g+R_{\Omega}+R_x)$\\
    欧姆表量程刻度&$R_x=\frac{E}{i}-\frac{E}{I_g}=E\frac{I_g-i}{iIg}$\\
    &$R_x=R_{\text{内}}(\frac{I_g}{i}-1)$\\
    欧姆表中值电阻&$R_{\text{中值电阻}}=\frac{E}{I_g}=R_{\text{内}}$\\
\end{tabularx}
\textbf{注意}\\
\begin{tabularx}{\columnwidth}{XX}
	\rowcolor{blue!30}
	\multicolumn{2}{l}{\parbox[c][2.5\ht\strutbox][c]{0.9\columnwidth}{\Large 用电器功率和电源功率}} \\
    电功率定义式&$P=UI$\\
    \multicolumn{2}{l}{\textbf{用电器功率}}\\
    纯电阻功率&$P=P_{\text{热}}=I^2R=\frac{U^2}{R}$\\
    非纯电阻总功率&$P=UI$\\
    非纯电阻热功率&$P=I^2R$\\
    非纯电阻输出功率&$P=UI-I^2R$\\
    \multicolumn{2}{l}{\textbf{电源功率}}\\
    电源总功率&$P=EI=\frac{E^2}{R+r}\le \frac{E^2}{r}$\\
    电源热功率&$P=I^2r=(\frac{E}{R+r})^2r\le \frac{E^2}{r}$\\
    电源输出功率&$P=UI=(E-Ir)I=(\frac{E}{r}-I)Ir$\\
    包含纯电阻电源输出功率&$P=(\frac{E}{R+r})^2R=\frac{E^2}{R+\frac{r^2}{R}+2r}\le \frac{E^2}{4r}$\\
\end{tabularx}
\textbf{注意}\\
\begin{tabularx}{\columnwidth}{XX}
	\rowcolor{blue!30}
	\multicolumn{2}{l}{\parbox[c][2.5\ht\strutbox][c]{0.9\columnwidth}{\Large 电路实验}} \\
    电流表外接&$R_x<\sqrt{R_VR_A}$\\
    通过试触法判断&$\frac{\Delta I}{I} < \frac{\Delta U}{U}$\\
    电流表内接&$R_x>\sqrt{R_VR_A}$\\
    通过试触法判断&$\frac{\Delta I}{I} > \frac{\Delta U}{U}$\\
    滑动变阻器限流解法&$R_{\text{滑动变阻器}}$略大于或者为3-5倍 $R_x$\\
    滑动变阻器分压解法&$R_{\text{滑动变阻器}}<R_x$\\
    &U需要从0开始\\
\end{tabularx}
\textbf{注意}\\
\begin{tabularx}{\columnwidth}{XX}
	\rowcolor{blue!30}
	\multicolumn{2}{l}{\parbox[c][2.5\ht\strutbox][c]{0.9\columnwidth}{\Large 动量}} \\
    
\end{tabularx}
\textbf{注意}\\
\begin{tabularx}{\columnwidth}{XX}
	\rowcolor{blue!30}
	\multicolumn{2}{l}{\parbox[c][2.5\ht\strutbox][c]{0.9\columnwidth}{\Large 磁场}} \\
    
\end{tabularx}
\textbf{注意}\\
\end{document}